\section{The two remaining compatibility conditions on the original member}

% BEGIN T86C1
\paragraph{T86C1. Recovering the parabolic from the character on a nonzero form.}
Fix a selected member, a coefficient index and a designated place above
one of the two type primes $g$. Put $n=r+t$, $k=\mathbb F_g$ and let
$U\subset\mathrm{GL}_n(k)$ consist of the block upper triangular
matrices having upper unitriangular first block of size $r$ and
arbitrary invertible second block. The principal and geometric frames
on this space give reductions related by $a\mapsto bab^{-1}$ for some
$b\in\mathrm{GL}_n(k)$. Indeed the standard frame in the member's
adelic realization identifies the action by conjugation over the local
field and carries the maximal compact into the stabilizer of the
lattice in the geometric frame. Every local integral invertible
matrix occurs, by inserting it at the split place with the
inverse-transpose partner and identity elsewhere. In the elementary
divisors of the frame change all exponents must be equal: otherwise
some elementary transvection would have a nonintegral conjugate.
Rescale the change of frame by a uniformizer power and reduce it.

If the row space is nonzero, $b$ preserves the first $r$ coordinates
as a subspace. To see this, take a nonzero realized form and a
linear functional on its value space detecting one of its values.
The resulting scalar function transforms through a character for
the inverse images of $U$ in both frames. In the principal frame
this is the specified level character, nontrivial on each simple
upper transvection in the first block by its additive character
on the sum of the superdiagonal. In the geometric frame this follows
by writing the same form using the member's adelic comparison:
the defining geometric level character acts by scalars on this
level subgroup. Matrices reducing to one fix the function: their
upper superdiagonal entries vanish in $k$, and the lower
determinant character on units is tame. These assertions only
use adeles supported at the tested place. Thus the stabilizer of
the line of this scalar function has a character on its reduction
$H$ containing $U$ and $bUb^{-1}$, nontrivial on each conjugated
upper simple root group $bU_{j,j+1}b^{-1}$, $1\leq j<r$.

Here is a short proof of the finite-group step. Write $b=b_1wb_2$
by Bruhat decomposition with $b_1,b_2$ upper triangular and $w$
a permutation. Upper triangular matrices normalize $U$. Conjugating
$H$ by $b_1^{-1}$ gives a group containing $U$ and $wUw^{-1}$.
The character on the latter is still nontrivial on each
$U_{w(j),w(j+1)}$: on the top block, upper triangular conjugation
scales each simple superdiagonal on the abelianization of the
upper unitriangular group. Its extra upper-right block contributes
no character (conjugation by a scalar on the lower block scales
that additive block by a scalar different from one). The bottom
diagonal block of a conjugated top transvection is the identity.
If either endpoint of such an edge were in the last $t$ coordinates,
a diagonal matrix in $U$ could scale this root group by $d\ne0,1$
by scaling just that lower coordinate. A character invariant under
this conjugation kills the whole group, since $(d-1)k=k$.
The type primes give $|k|>2$ and $r\ge2$, so these edges meet
every one of the first $r$ coordinates. The permutation preserves
their set, proving the assertion about $b$.
% END T86C1

% BEGIN T86C2
\paragraph{T86C2. The complex central action is the specified sum of right translations.}
At each designated place lift the first-block matrix of the finite
type group $G$ to an invertible integral matrix and adjoin the identity
lower block. Use the split unitary partner and take identity outside
the designated places. Denote these adelic lifts by $L(a)$, $a\in G$.
At auxiliary places their reduction is one. For a nonzero row, write
$A$ for the top block of the change $b$ in T86C1. In the geometric
frame $L(a)$ is block upper triangular with diagonal blocks
$AaA^{-1}$ and $1$. Conjugating the input $a$ by $A^{-1}$ at
each place gives the correct top coordinate simultaneously by a single
inner automorphism of $G$. The compensating level character in the
geometric action is one: the lower determinant blocks are identities,
the auxiliary reductions are one, and the top character of the
geometric flag level in this calculation is trivial.

For a rational conjugacy-class sum, the row's given integral-coordinate
comparison $\mathrm{AllT.cen}$ identifies the operator on rational
forms with its central level matrix. Under inclusion into ambient
coordinates that matrix acts as the sum of the corresponding finite
translations: the class sum commutes with the idempotent defining
the level. Under evaluation of an ambient coordinate form, the
translation is right translation by any lift multiplied by the inverse
level character of the lift, by the equivariance of the defining
function. The adelic comparison of the actual member identifies
these evaluated forms in the two frames. Use the lifts above and
reindex the class sum by the one inner automorphism. Thus for any
$z\in Z(\mathbb C[G])$ and form $F$ in this row,
\[
 \mathrm{real}(\mathrm{cenMt}(z)F)=
 \sum_{a\in G}z_a R(L(a))\,\mathrm{real}(F).
\]
Indeed the integral class sums span the complex center and the
rational coordinate basis spans the complex forms. For a zero row
the equality holds with the original lifts as well. This proves
$\mathrm{CenLoc}$ on the original localized module at every index.
The type space of each representative is the inverse image of its
isotypic sum under $\mathrm{real}$. Each right translation
preserves that sum. Thus every central operator preserves each type
space. The convention $\mathrm{onMt}=0$ away from the center
also gives the full stability input in M84T3. Together with
M84T3--4 and M84T7, this produces the integral action of the
center, its Hecke commutation and the scalar action on the whole
rational Hecke eigenspace at the original index.
% END T86C2

% BEGIN T86C3
\paragraph{T86C3. A local block vector and its two finite series.}
Fix an original representative $P$, including one with no Hecke
point, and a designated place $w$. Write $K=\mathrm{GL}_n(\mathbb
Z_g)$, $K_1=\ker(K\to\mathrm{GL}_n(k))$. Let $\pi$ be the
actual local factor of $P$ in the preferred split coordinate.
For $j=j(i)-1$ in zero-based notation let $\xi$ be the local
component of the tuple character $\mu_j^0\xi_{\rm tw}$, and
$\theta$ its character on $k^\times$. At either type prime
$\xi_{\rm tw}$ is unramified: the auxiliary character is supported
at an auxiliary prime distinct from both type primes. When
$\ell_i=g$ the remaining twisting character is trivial on units
above $\ell_i$; when $\ell_i\ne g$ the place belongs to $S_i$,
disjoint from its permitted ramification set $Q_i$. Hence
$\theta$ reads $\mu_j^0$ there. With the original residue
generator $h_1$, define $c\in k^\times$ by
$\theta(h_1)=\mathrm{Teich}(c)$ in the finite-series convention.
The tame-generator pin identifies $h_1$ with the residue of
$s(g^{1/(g-1)})/g^{1/(g-1)}$. A global integer unit at $w$
lifting this residue passes the exponent-one tame test as in
M84T10; the member's character test gives
$c\in\{\chi_J,\chi_J^{-1}\}$. The rank-one case of G85M2
with the same reciprocity convention identifies
$X-\mathrm{Teich}(c)$ with the inertia polynomial of
$\mathrm{locG}(\mu_C(i),w)$ at the named generator. Here the
actual character corresponds to the tuple by the member's
character identification; any unramified normalization is one on
inertia.

If a central top-block element sends a vector in $\mathrm{typeV}(i,P)$
to a nonzero vector, then there is $v\in\pi^{K_1}$ on which the
corresponding top action is nonzero, fixed by the parabolic
unipotent radical and transforming by $\theta\circ\det$ on the
lower block. Indeed in the level test the simple superdiagonal
character is trivial for lifts of this radical and of the lower
block, leaving respectively $1$ and $\theta\circ\det$.
T86C2 identifies the localized central operator with the finite
sum of top-block translations. For a lift of a single-place matrix
the other components may be taken to be identities; changing a
lift by the principal congruence group fixes the vector. Project
the resulting nonzero translated sum to a simple copy detecting
it in the global isotypic space, and then to a copy of $\pi$ in
the restriction to the local factor. Such copies span: translates
away from the place of an embedded local factor span a nonzero
global invariant subspace of the global irreducible. A sum of
simple modules is semisimple and the projections commute with the
local action. Project the original vector by the same maps.

In the series convention of the row, e.g. $\mathrm{grp.ser}$,
the idempotent for $f$ is the sum of the primitive central idempotents
of the finite irreducibles with Green label $f$. Finite parabolic
induction multiplies the cuspidal-support polynomials, by transitivity
of induction and uniqueness of finite cuspidal support. In rank $t$
the determinant character has label $(X-c)^t$: it occurs in
induction from the $t$ equal torus characters $\theta$, by
Frobenius reciprocity. The full inertia polynomial of any series
detected in $\pi^{K_1}$ is its lifted root polynomial at this same
generator. Apply the proof of G85M2 with ambient rank $n=r+t$.
Its only rank restriction was on an individual support factor $d$;
here $d\le n_V$ and $g-1$ is divisible by $Q_S q d_{\rm host}$
with $16\mid Q_S$, $q,d_{\rm host}\ge n_V$. Its numerical guard
therefore still holds. The member's $\mathrm{au}$, $\mathrm{embAV\_loc}$
and Weil--Deligne reading give precisely $\mathrm{locV}(i,P,w)$
in this computation, as in D85G1. These arguments supply the system
and all clauses of $\mathrm{LocType}$ in M84T8.
% END T86C3

% BEGIN T86C4
\paragraph{T86C4. The fourth tie.}
On the original member $D=w.D(Z)$, $\mathrm{SerLoc}(D)$ holds.
In the notation of T86C3 a top projector for $f$ detected on a
parameter space is detected in a representative's type space,
by the defining span for that parameter. Its projected vector is
fixed by the top projector, by the lower determinant projector,
and by the unipotent radical. Finite parabolic induction therefore
puts it in the series of $f(X)(X-c)^t$. On the member's base-change
list the rank-one string contributes $(X-\mathrm{Teich}(c))^t$;
all real shifts and the $\eta_T^{-1}$ twist preserve inertia.
Cancellation in the polynomial ring gives the lifted root polynomial
of $f$ on the remaining parameter. M84T8 proves equality with the
native monic label, including multiplicities, using the faithful
Teichm\"uller reading. T86C3 supplies its former hypotheses at
each original place and index. This is the original test for every
$F\in\mathrm{MtPar}$ and every polynomial $f$.
% END T86C4

% BEGIN T86C5
\paragraph{T86C5. Monodromy on the repeated determinant character.}
Suppose an actual finite irreducible top type $\tau$ occurs on
the space for $P$ in T86C3, with squarefree label $f$ avoiding
both prescribed values $\chi_J^{\pm1}$. A vector of the preceding
block kind can be projected to the $\tau$-isotypic space with
nonzero result. The lower block commutes with this projector and
the upper block normalizes the radical. Thus the parabolic
radical-invariants of $\pi^{K_1}$, restricted to its Levi, contain
$\tau\boxtimes (\theta\circ\det)$.
We show that the monodromy on the entire $n$-dimensional
parameter vanishes. The $t$ repeated copies of the same inertia
eigenvalue require an argument in addition to squarefreeness of $f$.

Write $A,B$ for the ranks $r,t$ finite representations
$\tau,\theta\circ\det$, and write $A*B$ for finite
Harish-Chandra induction. The finite product $A*B$ is irreducible;
moreover $A*B\simeq A'*B'$ with $A',B'$ irreducible in the same
respective cuspidal series forces $A'=A$, $B'=B$ up to
isomorphism. For clarity, Mackey's formula and Frobenius
reciprocity compute the intertwining space between these two
inductions by the double cosets of the $(r,t)$ parabolic.
Its intersections with a conjugate parabolic give the two
parabolic restrictions in the respective Levi blocks: on vectors
one takes invariants under the two intersecting unipotent radicals,
using their averaging idempotents in characteristic zero. In a cell
exchanging $a>0$ dimensions the exchanged $\mathrm{GL}_a(k)$
factor would be a constituent of a parabolic restriction from
each of the two series. Its cuspidal support would be a nonempty
submultiset of each original support, by transitivity and finite
Frobenius reciprocity. One support consists of $t$ copies of
$X-c$, while the other has no such root. Hence only the identity
cell contributes, giving
$\mathrm{Hom}(A,A')\otimes\mathrm{Hom}(B,B')$. Taking self
intertwiners and using finite semisimplicity proves irreducibility.
By the block vector and Frobenius reciprocity, $A*B$ occurs in
$\pi^{K_1}$.

Take the Zelevinsky multisegment of $\pi$. Its cuspidal factors
are depth zero, and their Green polynomials, including
multiplicities, have product $f(X)(X-c)^t$, by the support
calculation in G85M2 and T86C3. A factor with Green polynomial
$X-c$ is a rank-one character restricting to $\theta$ on units.
Indeed its Green polynomial has degree the cuspidal rank and
the degree-one Green correspondence reads the unit character.
Conversely every such character contributes $X-c$. Group the
segments on this unramified-twist line together, and put the
others in the first group. Segments in different groups are
unlinked. The multisegment classification and the irreducibility
of products on distinct cuspidal lines give
$\pi\simeq\pi_1\times\pi_2$ (normalized induction), with block
ranks $r,t$ and the second block on that line. Taking $K_1$
invariants gives finite parabolic induction of
$\pi_1^{K_{1,r}}\boxtimes\pi_2^{K_{1,t}}$ by Iwasawa
decomposition as in G85M2. All finite constituents of these two
spaces have respective labels $f$ and $(X-c)^t$ by the same
support argument. Decompose into finite irreducibles and apply
the finite intertwining calculation to the constituent $A*B$.
It gives $A\subset\pi_1^{K_{1,r}}$ and
$B\subset\pi_2^{K_{1,t}}$. In particular
$\pi_2\otimes(\xi^{-1}\circ\det)$ has a hyperspecial fixed
vector. Its spherical Weil--Deligne parameter has zero monodromy
by the unramified correspondence. Twisting back by a character
preserves this vanishing.

For $\pi_1$ the named inertia operator has characteristic
polynomial equal to the lift of $f$, with distinct roots.
Monodromy commutes with inertia, preserves every one-dimensional
eigenspace, and is nilpotent. Hence it is zero. The segment law
gives the direct sum of the two parameters for the irreducible
product. Thus $\mathrm{mono0}$ holds for the actual local
parameter, and its native label is $f(X)(X-c)^t$ by T86C3 and
the full-degree tame inertia reading of M84T9. Tameness here
follows either from the depth law applied to the $K_1$-fixed vector
as in D85G1 or from the individual depth-zero support factors.
At the conjugate place the paired split coordinates are exchanged
by inverse transpose. For irreducible admissible
$\mathrm{GL}_n$ this gives the contragredient by the
Gelfand--Kazhdan theorem. The duality and field-transport laws
of the correspondence apply: every support factor has
principal-congruence invariants as in D85G1. They give the dual
parameter under the conjugate field identification. Its monodromy
is the negative transpose of zero. M84T9 transports the designated
generator, so its label has degree $n$ and the inverses of the
roots of $f(X)(X-c)^t$, with multiplicity. The same $c$ works
in both conclusions, exactly as required in $\mathrm{admDZ}$.
% END T86C5

% BEGIN T86C6
\paragraph{T86C6. The integral witness, its parameter map and the second tie.}
Fix the local correspondence and the eventual design range of
M84T5. Choose for $\mathrm{lz}$ the finite irreducible
isomorphism classes, their ordinary character idempotents and
their Green labels in all ranks, with representations up to
isomorphism, ordinary occurrence, and finite parabolic induction.
In characteristic zero the finite group algebra is a product of
matrix algebras; these are all the primitive central idempotents,
orthogonal with sum one. The Green labels are monic of the
specified ranks with nonzero constant terms, and the member's
rank-$r$ series idempotents are the sums over their indicated
labels by T86C3. Take $\mathrm{kty}$ to be exactly the
irreducibles whose idempotents act nontrivially on some vector of
the given $\mathrm{typeV}$. Then T86C5 is the entire
$\mathrm{admDZ}$ assertion for this choice.

We supply the parameter-surjectivity field for the same witness.
If $\Sigma$ occurs at the original point $x$, the definition
supplies a representative $P$ over $x$ for which
$(\Sigma\sqcup\{(1,\mu_C(i),t)\})\otimes\eta_T^{-1}$
is a base change. Its parts have positive ranks and lengths,
are conjugate self-dual and are pairwise distinct as triples.
Twisting back is inverse on actual cuspidal representations
since the two characters multiply to one; this identity reaches
the original labels by the member's cuspidal identification
(equivalently by strong multiplicity one on the split-good
readings of the same global representations). Twisting by a
conjugate self-dual character preserves conjugate self-duality.
Deleting the final part and this common twist thus gives the
required positivity, self-duality and distinctness for $\Sigma$.
Its rank is $n-t=r$. At every designated type place the local
parameter of $P$ is tame: the local $K_1$-fixed-vector proof
in D85G1 uses the selected level at that place for every $i$,
irrespective of whether $g=\ell_i$. In the base-change sum
$\mathrm{isoLoc}(\Sigma,w)$ becomes a direct summand after an
unramified twist. Wild inertia therefore acts trivially on this
summand. In particular its monic label has degree $r$ at every
designated place by the wild-invariant label description in
M84T9. These are all the defining tests of the member's
concrete parameter set $\mathrm{ParC}$. Taking the inverse image
of this very list under the parameter identification
$\mathrm{par}.e_P$ gives one native parameter having the required
local values at \emph{every} place by its localization identity.
This proves $\mathrm{Par\_surj}$.

Finally take the roots of unity and coefficient-field series
coefficients of M84T5, the inertia characters and both polynomial
identities of M84T9, and the rational and integral actions and
projected eigenspaces of T86C2. The $\mathrm{EigSpan}$ assertion
is M84T3 on the same integral action. These choices give one
$Q\in\mathrm{Pend}(D)$ with both $\mathrm{admDZ}(Q)$ and
$\mathrm{EigSpan}(Q)$. In the eventual member quantifier the
finite list used in M84T5 was prescribed from the base before
the choice of $w,Z$. All later constructions are on the given
member and require no enlargement depending on $i,P$, a type
or a tested polynomial. Together with M84T10, M84T2 and T86C4
this proves all four original ties before making a good cut.
% END T86C6
